\chapter{\IfLanguageName{dutch}{Performantieanalyse en optimalisatie}{Performance Analysis and Optimization}}%
\label{ch:performantieanalyse-en-optimalisatie}
In dit hoofdstuk wordt de performantie van de ontwikkelde realtime monitoringapplicatie onderzocht en geoptimaliseerd. Hierbij wordt nagegaan welke factoren invloed hebben op de snelheid en efficiëntie van het systeem bij het verwerken en visualiseren van realtime voertuigdata.
De analyse richt zich voornamelijk op de invloed van de updatefrequentie en de hoeveelheid realtime data op de werking van de applicatie. Verschillende instellingen worden getest en met elkaar vergeleken om te bepalen welke configuratie een goede balans biedt tussen een vlotte realtime weergave en een beperkte systeembelasting.
Op basis van de resultaten worden waar nodig optimalisaties uitgevoerd. Het doel is om een stabiele en performante configuratie te bepalen die geschikt is voor de verdere toepassing van het monitoringdashboard.

\section{Performantiecriteria}
Hier bespreek je:
- welke aspecten je als performantie beschouwt
- latency
- responstijd
- updatefrequentie
- CPU-belasting
- geheugengebruik
- netwerkverkeer
- vloeiendheid van het dashboard
- waarom deze criteria relevant zijn voor realtime voertuigmonitoring

\section{Invloed van de updatefrequentie}
Hier onderzoek je:
- verschillende updatefrequenties
- bijvoorbeeld 100 ms, 250 ms, 500 ms en 1000 ms
- hoeveel updates per seconde dit oplevert
- invloed op de realtime weergave
- invloed op CPU-belasting
- invloed op netwerkverkeer
- invloed op de hoeveelheid verwerkte data

\section{Verwerking van realtime data}
Hier bespreek je:
- hoe nieuwe voertuigdata wordt verwerkt
- hoeveel data gelijktijdig wordt verwerkt
- hoe data wordt opgeslagen of bijgehouden
- hoe oude of overbodige data wordt behandeld
- eventuele technieken om onnodige verwerking te beperken
- hoe de backend en frontend met continue datastromen omgaan

\section{Optimalisatie}
Hier bespreek je:
- welke optimalisaties je hebt toegepast
- waarom deze optimalisaties nodig waren
- welke instellingen zijn aangepast
- eventuele optimalisatie van de updatefrequentie
- beperken van onnodige requests of berichten
- beperken van onnodige frontend-updates
- eventuele caching of buffering
- het verschil vóór en na de optimalisaties

\section{Evaluatie van de optimale configuratie}
Hier bepaal je uiteindelijk:
- welke updatefrequentie het meest geschikt is
- welke configuratie de beste balans biedt
- welke impact de optimalisaties hebben gehad
- hoe de uiteindelijke configuratie presteert
- waarom deze configuratie wordt gebruikt in de definitieve applicatie